% !TEX root = ../enac_project_fr.tex
\chapter{Linéarisation et méthodes d'estimation}
\label{ann:estimation}
\label{annexB}

Cette annexe regroupe les développements de linéarisation, de moindres carrés
pondérés et de filtrage de Kalman précédemment présentés dans le corps du
rapport. Les modèles physiques sont décrits au chapitre~\ref{chap:modelisation}.
Dans le développement géométrique initial, les corrections sont supposées
évaluées et les ambiguïtés sont notées $B_r^s$ pour alléger l'écriture.
En PPP iono-free, cette notation désigne $B_{r,\mathrm{IF}}^s$.
Les paramètres atmosphériques estimés sont ensuite ajoutés au modèle.

\section{Linéarisation du modèle}
\label{ann:resolution}

Pour préparer la résolution, on introduit deux fonctions reliant les observations aux paramètres à estimer. On note :
$f_s$ pour le code et $g_s$ pour la phase. Les
modèles corrigés s'écrivent

\begin{align}
\overline P_r^s
&=f_s(x_r,y_r,z_r,c\delta t_r)+\varepsilon_{P,r}^s,
\\
\overline L_r^s
&=g_s(x_r,y_r,z_r,c\delta t_r,B_r^s)+\varepsilon_{L,r}^s.
\label{ann:eq:explicit_observation_functions}
\end{align}

avec :

\begin{itemize}
    \item $\overline P_r^s$ et $\overline L_r^s$ les observations corrigées.
    \item $\varepsilon_{P,r}^s$ et $\varepsilon_{L,r}^s$ leurs erreurs
    résiduelles.
    \item $f_s$ et $g_s$ les fonctions d'observation dont les expressions
    sont précisées ci-après.
\end{itemize}


Pour expliciter la dépendance de ces fonctions aux paramètres, on note
$(x^s,y^s,z^s)$ les coordonnées du satellite $s$ et
$(x_r,y_r,z_r)$ celles du récepteur. Les fonctions d'observation de code et de
phase sont définies par

\begin{align}
f_s(x_r,y_r,z_r,c\delta t_r)
&=
\sqrt{(x_r-x^s)^2+(y_r-y^s)^2+(z_r-z^s)^2}
+c\delta t_r,
\label{ann:eq:code_observation_function}
\\
g_s(x_r,y_r,z_r,c\delta t_r,B_r^s)
&=
f_s(x_r,y_r,z_r,c\delta t_r)+B_r^s.
\label{ann:eq:phase_observation_function}
\end{align}

avec :

\begin{itemize}
    \item $f_s$ la fonction reliant la pseudo-distance corrigée à la position et au biais d'horloge du récepteur.
    \item $g_s$ la fonction reliant la phase corrigée à la position, au biais d'horloge et à l'ambiguïté métrique associée au satellite $s$.
    \item $(x^s,y^s,z^s)$ les coordonnées connues du satellite dans le même
    repère que le récepteur.
    \item $(x_r,y_r,z_r)$ les coordonnées inconnues du récepteur.
    \item $B_r^s$ l'ambiguïté métrique de la mesure de phase.
\end{itemize}


Les fonctions $f_s$ et $g_s$ dépendent de façon non linéaire des
coordonnées inconnues du récepteur et ne peuvent pas être résolues
comme un système linéaire. La démarche consiste à partir d'un état approché,
puis à calculer une correction au moyen d'un développement limité au premier
ordre. Les paramètres à estimer sont d'abord regroupés dans le vecteur d'état

\begin{align}
\mathbf x
&=
\begin{bmatrix}
\mathbf r_r^{\mathsf T} & c\delta t_r & B_r^1 & \cdots & B_r^n
\end{bmatrix}^{\mathsf T}.
\label{ann:eq:static_state_vector}
\end{align}

avec :

\begin{itemize}
    \item $\mathbf x$ le vecteur des paramètres estimés.
    \item $\mathbf r_r=[x_r\ y_r\ z_r]^{\mathsf T}$ la position du
    récepteur, en mètres.
    \item $c\delta t_r$ le biais d'horloge du récepteur, en mètres.
    \item $B_r^s$ l'ambiguïté associée au satellite $s$, en mètres.
\end{itemize}


On choisit alors une valeur approchée $\mathbf x_0$ de ce vecteur. Au voisinage
de ce point de linéarisation, le modèle non linéaire
$\mathbf h(\mathbf x)$ est approché au premier ordre par

\begin{align}
\mathbf y-\mathbf h(\mathbf x_0)
&\simeq
\mathbf H_0\,\Delta\mathbf x+\boldsymbol\varepsilon,
\\
\mathbf H_0
&=
\left.
\frac{\partial\mathbf h}{\partial\mathbf x}
\right|_{\mathbf x_0},
\\
\Delta\mathbf x
&=
\mathbf x-\mathbf x_0.
\label{ann:eq:linearized_model}
\end{align}

avec :

\begin{itemize}
    \item $\mathbf y$ le vecteur des observations corrigées, en mètres.
    \item $\mathbf h(\mathbf x_0)$ le vecteur des observations calculées à
    partir de l'état approché.
    \item $\mathbf H_0$ la matrice jacobienne du modèle.
    \item $\Delta\mathbf x$ la correction apportée à l'état.
\end{itemize}

En écriture scalaire, l'état approché autour duquel les fonctions sont
linéarisées est

\begin{align}
\mathbf x_0
&=
\begin{bmatrix}
x_{r,0} & y_{r,0} & z_{r,0} & c\delta t_{r,0} &
B_{r,0}^{1} & \cdots & B_{r,0}^{n}
\end{bmatrix}^{\mathsf T}.
\label{ann:eq:approximate_state}
\end{align}

avec :

\begin{itemize}
    \item $(x_{r,0},y_{r,0},z_{r,0})$ la position approchée du récepteur.
    \item $c\delta t_{r,0}$ le biais d'horloge approché.
    \item $B_{r,0}^{s}$ l'ambiguïté métrique approchée du satellite $s$.
\end{itemize}

La distance géométrique calculée à cette position vaut

\begin{align}
\rho_{r,0}^s
&=
\sqrt{(x_{r,0}-x^s)^2+(y_{r,0}-y^s)^2+(z_{r,0}-z^s)^2}.
\label{ann:eq:geometric_range_at_initial_state}
\end{align}

avec :

\begin{itemize}
    \item $\rho_{r,0}^s$ la distance calculée entre le satellite $s$ et la
    position approchée du récepteur.
    \item $(x^s,y^s,z^s)$ les coordonnées du satellite.
    \item $(x_{r,0},y_{r,0},z_{r,0})$ les coordonnées approchées du
    récepteur.
\end{itemize}

Les dérivées des deux fonctions par rapport aux coordonnées et au biais
d'horloge sont identiques. La fonction de phase possède en plus une dérivée
par rapport à son ambiguïté :

\begin{align}
\left.\frac{\partial f_s}{\partial x_r}\right|_{\mathbf x_0}
&=
\left.\frac{\partial g_s}{\partial x_r}\right|_{\mathbf x_0}
=\frac{x_{r,0}-x^s}{\rho_{r,0}^s},
\\
\left.\frac{\partial f_s}{\partial y_r}\right|_{\mathbf x_0}
&=
\left.\frac{\partial g_s}{\partial y_r}\right|_{\mathbf x_0}
=\frac{y_{r,0}-y^s}{\rho_{r,0}^s},
\\
\left.\frac{\partial f_s}{\partial z_r}\right|_{\mathbf x_0}
&=
\left.\frac{\partial g_s}{\partial z_r}\right|_{\mathbf x_0}
=\frac{z_{r,0}-z^s}{\rho_{r,0}^s},
\\
\frac{\partial f_s}{\partial(c\delta t_r)}
&=
\frac{\partial g_s}{\partial(c\delta t_r)}=1,
\\
\frac{\partial g_s}{\partial B_r^s}
&=1.
\label{ann:eq:component_observation_derivatives}
\end{align}


En introduisant les corrections suivantes :

\begin{align}
\Delta x_r&=x_r-x_{r,0},
&
\Delta y_r&=y_r-y_{r,0},
&
\Delta z_r&=z_r-z_{r,0},
\\
\Delta b_r&=c\delta t_r-c\delta t_{r,0},
&
\Delta B_r^s&=B_r^s-B_{r,0}^s.
\label{ann:eq:state_component_corrections}
\end{align}

Les deux observations du satellite $s$ sont linéarisées sous la forme

\begin{align}
\overline P_r^s-f_s(\mathbf x_0)
&\simeq
\frac{x_{r,0}-x^s}{\rho_{r,0}^s}\Delta x_r
+\frac{y_{r,0}-y^s}{\rho_{r,0}^s}\Delta y_r
+\frac{z_{r,0}-z^s}{\rho_{r,0}^s}\Delta z_r
+\Delta b_r+\varepsilon_{P,r}^s,
\\
\overline L_r^s-g_s(\mathbf x_0)
&\simeq
\frac{x_{r,0}-x^s}{\rho_{r,0}^s}\Delta x_r
+\frac{y_{r,0}-y^s}{\rho_{r,0}^s}\Delta y_r
+\frac{z_{r,0}-z^s}{\rho_{r,0}^s}\Delta z_r
+\Delta b_r+\Delta B_r^s+\varepsilon_{L,r}^s.
\label{ann:eq:component_linearized_observations}
\end{align}

En empilant d'abord les $n$ observations de code, puis les $n$ observations de
phase, le système linéarisé s'écrit :

\begin{align}
\Delta\mathbf y
&=
\mathbf H_0\Delta\mathbf x+\boldsymbol\varepsilon,
\label{ann:eq:stacked_linearized_system}
\\
\Delta\mathbf y
&=
\begin{bmatrix}
\overline P_r^1-f_1(\mathbf x_0)\\
\vdots\\
\overline P_r^n-f_n(\mathbf x_0)\\
\overline L_r^1-g_1(\mathbf x_0)\\
\vdots\\
\overline L_r^n-g_n(\mathbf x_0)
\end{bmatrix},
&
\Delta\mathbf x
&=
\begin{bmatrix}
\Delta x_r\\
\Delta y_r\\
\Delta z_r\\
\Delta b_r\\
\Delta B_r^1\\
\vdots\\
\Delta B_r^n
\end{bmatrix},
&
\boldsymbol\varepsilon
&=
\begin{bmatrix}
\varepsilon_{P,r}^1\\
\vdots\\
\varepsilon_{P,r}^n\\
\varepsilon_{L,r}^1\\
\vdots\\
\varepsilon_{L,r}^n
\end{bmatrix}.
\label{ann:eq:stacked_vectors}
\end{align}


La structure de la jacobienne est la suivante :

\begin{align}
\mathbf H_0
&=
\begin{bmatrix}
\mathbf G & \mathbf 0_{n\times n}\\
\mathbf G & \mathbf I_n
\end{bmatrix},
&
\mathbf G
&=
\begin{bmatrix}
\dfrac{x_{r,0}-x^1}{\rho_{r,0}^1} &
\dfrac{y_{r,0}-y^1}{\rho_{r,0}^1} &
\dfrac{z_{r,0}-z^1}{\rho_{r,0}^1} & 1\\
\vdots & \vdots & \vdots & \vdots\\
\dfrac{x_{r,0}-x^n}{\rho_{r,0}^n} &
\dfrac{y_{r,0}-y^n}{\rho_{r,0}^n} &
\dfrac{z_{r,0}-z^n}{\rho_{r,0}^n} & 1
\end{bmatrix}.
\label{ann:eq:gnss_design_matrix}
\end{align}

avec :

\begin{itemize}
    \item $\mathbf G$ la matrice géométrique reliant les observations aux
    trois coordonnées du récepteur et à son biais d'horloge.
    \item $\mathbf 0_{n\times n}$ la matrice nulle de taille $n\times n$.
    \item $\mathbf I_n$ la matrice identité de taille $n\times n$.
\end{itemize}

Plus la position initiale est proche de la solution, plus la linéarisation est valable. Une seule correction peut donc être insuffisante si la
position initiale est éloignée de la solution.

Le chapitre~\ref{chap:modelisation} décrit les corrections
appliquées avant la résolution. Une fois les observations corrigées et le
système linéarisé sous la forme de l'équation~\eqref{ann:eq:stacked_linearized_system},
les paramètres peuvent être estimés soit indépendamment à chaque époque par
moindres carrés pondérés, soit récursivement par filtre de Kalman étendu. Ces
deux méthodes sont présentées respectivement dans les
sections~\ref{ann:sec:wls} et~\ref{ann:Kalman}.

\section{Moindres carrés pondérés}
\label{ann:sec:wls}

La section~\ref{ann:resolution} a établi le système
linéarisé~\eqref{ann:eq:stacked_linearized_system} à partir des observations
corrigées. Les moindres carrés pondérés (ou \textit{Weighted Least Squares}, WLS) recherchent la correction qui rend les résidus
entre observations et valeurs calculées aussi faibles que possible, tout en
tenant compte de la précision attribuée à chaque observation. Le critère
minimisé est $J=\mathbf e^{\mathsf T}\mathbf W\mathbf e$, où
$\mathbf e=\Delta\mathbf y-\mathbf H_0\Delta\mathbf x$ est le vecteur des
résidus linéarisés. En notant
$\Delta\mathbf y=\mathbf y-\mathbf h(\mathbf x_0)$, l'annulation de la dérivée
de ce critère conduit à

\begin{align}
\widehat{\Delta\mathbf x}
&=
\left(\mathbf H_0^{\mathsf T}\mathbf W\mathbf H_0\right)^{-1}
\mathbf H_0^{\mathsf T}\mathbf W\Delta\mathbf y,
\\
\mathbf W
&=
\mathbf R^{-1}.
\label{ann:eq:wls_solution}
\end{align}

avec :

\begin{itemize}
    \item $\widehat{\Delta\mathbf x}$ la correction d'état estimée.
    \item $\mathbf W$ la matrice de pondération.
    \item $\mathbf R$ la covariance des observations.
    \item $\mathbf H_0$ la jacobienne évaluée à l'état courant.
\end{itemize}

Le choix $\mathbf W=\mathbf R^{-1}$ donne davantage de poids aux observations
les plus précises. Par exemple, une mesure dont la variance est élevée contribue
moins à la solution qu'une mesure de faible variance. L'inversion de la matrice
normale suppose que les paramètres soient observables avec la géométrie et le
jeu de mesures disponibles.

L'état est mis à jour itérativement selon

\begin{align}
\mathbf x_{j+1}
&=
\mathbf x_j+\widehat{\Delta\mathbf x}_j.
\label{ann:eq:wls_iteration}
\end{align}

avec :

\begin{itemize}
    \item $j$ l'indice d'itération.
    \item $\mathbf x_j$ l'état avant mise à jour.
    \item $\widehat{\Delta\mathbf x}_j$ la correction calculée à l'itération
    $j$.
\end{itemize}

Après chaque mise à jour, les distances et la jacobienne sont recalculées au
nouvel état. Les paramètres pouvant avoir des unités et des ordres de grandeur
différents, le critère d'arrêt porte sur la correction mise à l'échelle :

\begin{align}
\eta_j
&=
\max_p\left|
\frac{\widehat{\Delta x}_{j,p}}{s_p}
\right|,
&
\eta_j<\eta_{\mathrm{WLS}}
&\quad\Longrightarrow\quad
\text{convergence}.
\label{ann:eq:wls_stopping_criterion}
\end{align}

avec :

\begin{itemize}
    \item $\widehat{\Delta x}_{j,p}$ la correction du paramètre $p$ à
    l'itération $j$.
    \item $s_p$ l'échelle associée à ce paramètre.
    \item $\eta_j$ la plus grande correction normalisée.
    \item $\eta_{\mathrm{WLS}}$ le seuil de convergence.
\end{itemize}

Une valeur type est $\eta_{\mathrm{WLS}}=10^{-3}$. La résolution est répétée
jusqu'à satisfaction de ce critère ou jusqu'au nombre maximal d'itérations,
fixé ici à $j_{\max}=50$. Cette méthode produit une
solution complète à partir des seules observations de l'époque considérée.
Elle est donc adaptée à une résolution indépendante époque par époque, mais
n'exploite pas le fait que la position, les retards atmosphériques ou les
ambiguïtés évoluent généralement de manière continue. Cette limitation motive
l'introduction d'un modèle temporel.

\section{Filtre de Kalman étendu }
\label{ann:Kalman}

Le filtre de Kalman étendu (ou \emph{Extended Kalman Filter}, EKF) est une autre méthode de résolution des équations établies en section~\ref{ann:resolution}.
Il exploite la continuité temporelle : la solution
obtenue à une époque devient l'information a priori de l'époque suivante. Il
combine deux sources d'information. Le modèle d'évolution décrit la manière
dont l'état est susceptible de changer entre deux époques, tandis que le modèle
d'observation relie cet état aux mesures GNSS courantes. Leurs incertitudes
respectives déterminent le poids accordé à la prédiction et aux observations.

L'indice $k$
désigne désormais l'époque. Un exposant $-$ désignera une quantité prédite,
avant l'utilisation des observations de cette époque. L'absence d'exposant
désignera la quantité corrigée.

\subsection{Initialisation du filtre}
\label{ann:sec:ekf_initialization}

Le calcul doit commencer avec une estimation initiale
$\widehat{\mathbf x}_0$ et sa covariance $\mathbf P_0$. Chaque paramètre du
vecteur d'état doit donc recevoir une valeur initiale, même approximative. Dans
le cas d'une covariance initiale diagonale,

\begin{align}
\mathbf P_0
=
\operatorname{diag}
\left(
\sigma_{x,0}^2,
\sigma_{y,0}^2,
\sigma_{z,0}^2,
\sigma_{c\delta t,0}^2,
\sigma_{\mathrm{ISB},0}^2,
\sigma_{B_{r,0}^1}^2,\ldots,
\sigma_{B_{r,0}^n}^2,
\sigma_{ZWD,0}^2,\ldots
\right).
\label{ann:eq:ekf_initial_covariance}
\end{align}

avec :

\begin{itemize}
    \item $\sigma_{x,0}$, $\sigma_{y,0}$ et $\sigma_{z,0}$ les écarts-types
    initiaux des coordonnées.
    \item $\sigma_{c\delta t,0}$ et $\sigma_{\mathrm{ISB},0}$ ceux de
    l'horloge et, en multi-constellation, de l'ISB.
    \item $\sigma_{B_{r,0}^s}$ celui de l'ambiguïté métrique initiale
    $B_{r,0}^s$ du satellite $s$.
    \item $\sigma_{ZWD,0}$ celui du retard humide zénithal lorsque ce paramètre
    est estimé.
    \item les points de suspension désignent les autres paramètres éventuels, dont
    les gradients troposphériques.
\end{itemize}

La forme diagonale constitue une écriture simplifiée. Lorsque deux paramètres
sont initialement corrélés, leur covariance doit être placée dans les termes
hors diagonale de $\mathbf P_0$.

Les méthodes utilisées pour
obtenir les valeurs initiales propres au traitement GNSS sont présentées dans
la section~\ref{ppp:sec:ekf_gnss_features}.

\subsection{Modèle d'évolution}

Le modèle suivant exprime l'état courant à partir de l'état précédent :

\begin{align}
\mathbf x_k
&=
\mathbf F_k\mathbf x_{k-1}+\mathbf w_k,
\\
\mathbb E[\mathbf w_k]
&=\mathbf0,
&
\mathbb E[\mathbf w_k\mathbf w_k^{\mathsf T}]
&=\mathbf Q_k.
\label{ann:eq:kalman_prediction}
\end{align}

avec :

\begin{itemize}
    \item $\mathbf x_k$ l'état à l'époque $k$.
    \item $\mathbf F_k$ la matrice de transition entre les époques $k-1$ et
    $k$.
    \item $\mathbf w_k$ le bruit de processus.
    \item $\mathbf Q_k$ la covariance du bruit de processus.
\end{itemize}

La matrice $\mathbf F_k$ traduit l'évolution déterministe attendue. Pour un
récepteur statique, les coordonnées sont supposées constantes entre deux
époques. Leur matrice de transition est donc l'identité. Le bruit de processus
$\mathbf w_k$ représente les variations que ce modèle idéal ne décrit pas. Sa
covariance $\mathbf Q_k$ règle la liberté d'évolution accordée à chaque
paramètre : une faible variance impose une forte continuité, tandis qu'une
variance plus élevée permet une évolution plus rapide.

\subsection{Modèle d'observation}

À l'époque $k$, le vecteur $\mathbf y_k$ regroupe les observations disponibles.
La fonction $\mathbf h_k$ calcule les valeurs attendues à partir d'un état
donné et de la géométrie satellite-récepteur. Le modèle non linéaire des
observations est

\begin{align}
\mathbf y_k
&=
\mathbf h_k(\mathbf x_k)+\mathbf v_k,
\\
\mathbb E[\mathbf v_k]
&=\mathbf0,
&
\mathbb E[\mathbf v_k\mathbf v_k^{\mathsf T}]
&=\mathbf R_k.
\label{ann:eq:kalman_observation}
\end{align}

avec :

\begin{itemize}
    \item $\mathbf y_k$ le vecteur d'observations à l'époque $k$.
    \item $\mathbf h_k$ le modèle d'observation GNSS.
    \item $\mathbf v_k$ le bruit de mesure.
    \item $\mathbf R_k$ la covariance des observations.
\end{itemize}

La matrice $\mathbf R_k$ joue ici le même rôle que dans les moindres carrés
pondérés : elle décrit la précision et, le cas échéant, les corrélations des
mesures. Les bruits de processus et de mesure sont supposés indépendants.
Autour de l'état prédit $\widehat{\mathbf x}_k^-$, la relation linéarisée est

\begin{align}
\mathbf y_k-\mathbf h_k(\widehat{\mathbf x}_k^-)
&\simeq
\mathbf H_k
\left(\mathbf x_k-\widehat{\mathbf x}_k^-\right)
+\mathbf v_k,
\\
\mathbf H_k
&=
\left.
\frac{\partial\mathbf h_k}{\partial\mathbf x}
\right|_{\widehat{\mathbf x}_k^-}.
\label{ann:eq:ekf_linearization}
\end{align}

avec :

\begin{itemize}
    \item $\widehat{\mathbf x}_k^-$ l'état prédit avant traitement des
    observations.
    \item $\mathbf H_k$ la jacobienne du modèle à l'état prédit.
    \item le membre de gauche l'innovation non encore pondérée.
\end{itemize}

\subsection{Prédiction}

Avant de traiter les nouvelles mesures, le filtre propage l'estimation obtenue
à l'époque précédente. Il propage également sa covariance, qui quantifie
l'incertitude sur l'état. L'ajout de $\mathbf Q_k$ augmente cette incertitude
pour tenir compte des évolutions non parfaitement connues :

\begin{align}
\widehat{\mathbf x}_k^-
&=
\mathbf F_k\widehat{\mathbf x}_{k-1},
\\
\mathbf P_k^-
&=
\mathbf F_k\mathbf P_{k-1}\mathbf F_k^{\mathsf T}
+\mathbf Q_k.
\label{ann:eq:ekf_predict}
\end{align}

avec :

\begin{itemize}
    \item $\widehat{\mathbf x}_{k-1}$ l'état corrigé à l'époque précédente.
    \item $\mathbf P_{k-1}$ sa covariance.
    \item $\widehat{\mathbf x}_k^-$ et $\mathbf P_k^-$ l'état prédit et sa
    covariance.
\end{itemize}

\subsection{Mise à jour}

Lorsque les observations deviennent disponibles, le filtre les compare aux
valeurs calculées à partir de l'état prédit. Cette différence est appelée
\emph{innovation}. Une innovation nulle signifie que la prédiction explique
exactement la mesure. Une innovation importante peut provenir d'une erreur de
prédiction, du bruit de mesure ou d'une observation aberrante. L'innovation et
sa covariance sont

\begin{align}
\boldsymbol\nu_k
&=
\mathbf y_k-\mathbf h_k(\widehat{\mathbf x}_k^-),
\\
\mathbf S_k
&=
\mathbf H_k\mathbf P_k^-\mathbf H_k^{\mathsf T}+\mathbf R_k.
\label{ann:eq:ekf_innovation}
\end{align}

avec :

\begin{itemize}
    \item $\boldsymbol\nu_k$ le vecteur d'innovation.
    \item $\mathbf S_k$ la covariance de l'innovation.
    \item $\mathbf R_k$ la covariance du bruit de mesure.
\end{itemize}

La covariance $\mathbf S_k$ combine l'incertitude de l'état projetée dans
l'espace des observations et l'incertitude propre aux mesures. Le gain de
Kalman en déduit la part de l'innovation à appliquer à l'état :

\begin{align}
\mathbf K_k
&=
\mathbf P_k^-\mathbf H_k^{\mathsf T}\mathbf S_k^{-1},
\\
\widehat{\mathbf x}_k
&=
\widehat{\mathbf x}_k^-+\mathbf K_k\boldsymbol\nu_k,
\\
\mathbf P_k
&=
\left(\mathbf I-\mathbf K_k\mathbf H_k\right)\mathbf P_k^-.
\label{ann:eq:ekf_update}
\end{align}

avec :

\begin{itemize}
    \item $\mathbf K_k$ le gain de Kalman.
    \item $\widehat{\mathbf x}_k$ l'état après assimilation des observations.
    \item $\mathbf P_k$ la covariance de l'état corrigé.
    \item $\mathbf I$ la matrice identité de dimension adaptée.
    \item les autres termes sont définis dans les équations précédentes.
\end{itemize}

Si la prédiction est très incertaine devant les mesures, le gain accorde une
place importante à l'innovation. À l'inverse, des observations très bruitées
modifient peu l'état prédit. Le gain n'est pas choisi à l'avance. Il résulte
des covariances du modèle et des observations.

À l'issue de cette mise à jour, $\widehat{\mathbf x}_k$ et $\mathbf P_k$
résument toute l'information utilisée jusqu'à l'époque $k$. Ils servent de
point de départ à la prédiction suivante. Le cycle de prédiction et de mise à
jour est répété à chaque époque, ce qui affine progressivement les
paramètres constants tout en laissant évoluer ceux auxquels un bruit de
processus a été attribué.

Les mécanismes spécifiques au GNSS sont présentés dans la
section~\ref{ppp:kalman} du chapitre consacré au positionnement PPP.
